\section{A rational determinant and its lower bound}\label{sec:arithmetic}

Fix a rational \(\nu>3\) and fractions satisfying \eqref{eq:exceptional}, chosen
large enough that \(1<r_i<2\). Set \(w_i=\lceil\log q_i\rceil\)
and \(w_* =\min_iw_i\geq1\). Choose a rational \(F>2/v\), and use
\[
 T_i=\lceil Fw_i/v_0\rceil,\quad
 G_i(x)=\sum_{1\leq n<T_i}\frac{(-1)^{n+1}x^n}{n},\quad
 L_i=\lcm(1,\ldots,T_i).
\]
Since \(T_iv_0>w_i/v\),
substituting \(G_i\) for both logarithms preserves the interpolation
rank. The matrix, with all coefficient rows and all allowed monomial
columns, is rational and has full row rank for large divisible \(H\).
Choose a nonzero square minor \(\Delta_H\) using every row, and let its
size be \(M=M_H\).

Multiply column \((h,k,\alpha)\) by \(\prod_iq_i^{\alpha_i}\),
row \((j,\ell_1,\ell_2,\beta)\) by \(\prod_iq_i^{-\beta_i}\),
and every entry by
\(D_H=\prod_i L_i^{\lfloor H/w_i\rfloor}\).
The resulting entries are integers: apart from integral factors they
are coefficients of
\[
 2^{jh}3^{jk}\binom\alpha\beta(1+t)^h(1+s)^k
              \prod_i[q_iG_i(s)-p_iG_i(t)]^{\alpha_i-\beta_i}.
\]
Entries with some \(\beta_i>\alpha_i\) are zero.

Let \(\Lambda=4+\log4\), and put
\(\bar b_H=(MH)^{-1}\sum_{\rm rows}w\beta\leq v\).
The elementary bound \(\log\lcm(1,\ldots,n)\leq\Lambda n\)
suffices. For completeness, the ratio of the least common multiples
through \(2^k\) and \(2^{k-1}\) divides
\(\binom{2^k}{2^{k-1}}\): each prime power in \((2^{k-1},2^k]\)
contributes its underlying prime to that ratio and to the binomial
coefficient. Summing the logarithmic bound over dyadic intervals
and enlarging to the next power of two
gives the stronger constant \(4\log2\). The coarser constant
\(\Lambda\) will suffice. Nonvanishing of the integral determinant gives
\begin{equation}\label{eq:arith}
 \frac{\log|\Delta_H|}{MH}\geq-1+\bar b_H-E_{\rm ar},\qquad
 E_{\rm ar}=\frac{\Lambda Fm}{v_0}
          +\Lambda\sum_i\frac1{w_i}+\frac v{w_*}.
\end{equation}
Indeed the column logarithmic charge is at most \(MH\), the replacement
of \(\log q_i\) by \(w_i\) in the row credit loses at most
\(vMH/w_*\), and \(\log D_H/H\) is bounded by the first two terms
in \(E_{\rm ar}\).

Weighted-simplex lattice counts give, with all data fixed,
\begin{equation}\label{eq:counts}
 M=\frac{Kv^mH^{m+2}}{(m+2)!v_0^2\prod_iw_i}+O(H^{m+1}).
\end{equation}
This follows by rescaling the row simplex to the unit simplex, whose
volume is \(1/(m+2)!\).
